\documentclass[10pt,a4paper]{amsart}
\usepackage{amsmath,amssymb,amsthm,graphicx}
\RequirePackage{amssymb,amsmath,amsthm,bm}
\RequirePackage{thmtools}
\RequirePackage[capitalize]{cleveref}
\providecommand{\tightlist}{%
  \setlength{\itemsep}{0pt}\setlength{\parskip}{0pt}}

%%%%% varphi instead of phi
\renewcommand{\phi}{\varphi}

\mathchardef\mhyph="2D 
\providecommand{\A}{\mathcal{A}}
\providecommand{\C}{\mathcal{C}}
\providecommand{\L}{\mathcal{L}}
\providecommand{\B}{\mathcal{B}}
\providecommand{\M}{\mathcal{M}}
\renewcommand{\S}{\mathcal{S}}
\providecommand{\ol}[1]{\overline{#1}}
\providecommand{\mbf}[1]{\mathbf{#1}}
\providecommand{\mf}[1]{\mathfrak{#1}}
\providecommand{\mrm}[1]{\mathrm{#1}}
\providecommand{\mc}[1]{\mathcal{#1}}
\providecommand{\ra}{\rightarrow}
\providecommand{\Ra}{\Rightarrow}
\providecommand{\La}{\Leftarrow}
\providecommand{\lr}{\leftrightarrow}
\providecommand{\hf}{\mathbb{HF}}
\providecommand{\spec}{DgSp}
\providecommand{\bSigma}{\bm\Sigma}
\providecommand{\bDelta}{\bm\Delta}
\providecommand{\bPi}{\bm\Pi}
\providecommand{\Pinf}[1]{\Pi^{\mathrm{in}}_{#1}}
\providecommand{\Sinf}[1]{\Sigma^{\mathrm{in}}_{#1}}
\providecommand{\dSinf}[1]{d\mhyph\Sigma^{\mathrm{in}}_{#1}}
\providecommand{\Dinf}[1]{\Delta^{\mathrm{in}}_{#1}}
\providecommand{\bp}{\bar p}
\providecommand{\bq}{\bar q}
\providecommand{\ba}{\bar a}
\providecommand{\bb}{\bar b}
\providecommand{\bc}{\bar c}
\providecommand{\bd}{\bar d}
\providecommand{\bx}{\bar x}
\providecommand{\by}{\bar y}
\providecommand{\bz}{\bar z}
\providecommand{\LO}{\mathrm{LO}}
\declaretheorem[name={Theorem},numberwithin=section]{theorem}
\declaretheorem[name={Proposition}, sibling=theorem]{proposition}
\declaretheorem[name={Lemma}, sibling=theorem]{lemma}
\declaretheorem[name={Fact}, sibling=theorem]{fact}
\declaretheorem[name={Claim}, numberwithin=theorem]{claim}
\declaretheorem[name={Definition},numberwithin=section,sibling=theorem, style=definition]{definition}
\declaretheorem[name={Remark}, numbered=no,style=definition]{remark}
\declaretheorem[name={Corollary}, sibling=theorem]{corollary}
\declaretheorem[name={Question},sibling=theorem]{question}
\newcommand{\T}{{\mathcal{T}}}
\newcommand{\SR}{SR}
\newcommand{\pSR}{pSR}
\newcommand{\Q}{\mathbb{Q}}
\newcommand{\Z}{\mathbb{Z}}
\newcommand{\K}{\mathcal{K}}
\newcommand{\N}{\mathcal{N}}
\renewcommand{\L}{\mathcal{L}}
\providecommand{\Ginf}[1]{\Gamma^{\mathrm{in}}_{#1}}
\title{The natural numbers together with the rationals are two-universal for countable linear orders}
\author{David Gonzalez, Dino Rossegger}
\date{arxiv:2305.07126v1; CC BY 4.0}
\begin{document}
\maketitle
\noindent\textit{Editorial note.} The counted statement is the author environment at source0.tex lines 809-812 and its proof environment opens at line 813, the line immediately after; the pairing is by adjacency and the retained proof range 813-818 is the whole environment.\par\medskip
\noindent\textit{Editorial note.} The retained preamble is the author style file dino.sty (package lines, the theorem declaration support, the notation and tuple macros and the theorem environments) followed by the ordinal and rank macros of the manuscript itself, reproduced verbatim; this is the same recipe as the previously published card of this article, so no package outside it is requested and extra\_packages stays empty.\par\medskip
\noindent\textit{Editorial note.} Metadata provenance: source\_id and source\_version are the values recorded in this source licence evidence file (sources/scott-linear-orders-2305.07126v1/LICENSE\_EVIDENCE.json: arxiv:2305.07126v1 / 2305.07126v1), which the card points at directly; the title and author list are transcribed from the authors\_index and title keys of the source record of the pinned copy. The card carries no counter setting and imposes no original numbering, since this article has no label-number dump and its numbering cannot be reconstructed from the retained ranges alone.\par\medskip
\noindent\textit{Editorial note.} No context span and no external reference are needed: the retained statement and proof use no label or citation from outside their own ranges.\par\medskip
\noindent\textit{Editorial scope.} One original result of the pinned author manuscript, counted once: two-universality of the natural numbers together with the rationals for countable linear orders (author label lem:universal).\par\medskip
\noindent\textit{Allowed standard prerequisites (editorial).}\begin{itemize}
\item Two-universality. A class K is two-universal for the class of linear orders when every countable linear order stands in the rank two back-and-forth relation to some member of K.. It is what the lemma establishes for N together with Q.
\item Density of the rationals. Every interval of Q is infinite and has the order type of Q.. It supplies the rank one comparisons used in the proof.
\end{itemize}\medskip
\begin{lemma}\label{lem:universal}
  The class $K=\mathbb N\cup\{\Q\}$ is $2$-universal for $\LO$. In particular,
  $\{\Q\}$ is $2$-universal for infinite linear orderings.
\end{lemma}

\begin{proof}
  Let $\A\in \LO$. If $\A$ is finite then $\A\in K$. If
  $\A$ is infinite then we claim that $\A\leq_2\Q$. As any interval $I$ in
  $\Q$ is infinite we have that $I\leq_1\A_0$ for any interval $\A_0$ of $\A$.
  Thus, $\A\leq_2\Q$ and, hence, $K$ is $2$-universal.
\end{proof}

\end{document}
